\chapter{Chapter 3 - Enemy Reborn}
\label{chap:chapter3}

\setcounter{section}{-1}
\section{Running This Chapter}

In this chapter, the riots continue on the streets of Kholinar. \npcHighmarshal{} defends city gates and storms the monastery, while helpless parshmen are slaughtered on the streets by queen's orders.

In this chaos, the characters follow their own agenda - they must reclaim ancient artifacts of the Unmade from the Ardent Monastery.

\reasonBox{
	Starting from this chapter, this book assumes \npcTroublemaker{} might have died in the riots and she is referenced no more. If in your game she survived, it is up to you how to tell her story - she may die from the hands of newborn Singers during \fullref{sec:enemyReborn}, or she may live to meet the PCs again during \nameref{chap:chapter9}. Who knows, she might even bond an Enlightened spren (knowing her personality, probably a lightspren)
}

\subsection{Character Advancement}

The characters start this chapter at level 3. Advance them to level 4 at the end of this chapter.

\reasonBox{
	Depending on the player choices, this chapter may be encounter-rich, or finish quicker if the party remains focused to complete their objective as soon as possible.
	\\ \\
	If you feel the chapter is getting longer, consider advancing the party to level 4 as you trigger \fullref{sec:enemyReborn} events. If you do so, do not award them additional level at the end of the chapter, they should start \nameref{chap:chapter4} at level 4 no matter how many encounters did they have during the riots.
}

\newpage
The following events may affect their path to radiance:
\begin{itemize}
	\item Calming angry mob or leading them to riot against queen's herald
	\item Defending helpless parshmen or helping to execute them
	\item Fighting with \npcThrill{} nearby can expedite an Enlightened Windrunner
	\item Seeing ancient artifacts of the Unmade can expedite an Enlightened Dustbringer or Elsecaller
\end{itemize}



\subsection{Forging Alliances}

The following actions can help the characters to become allies with Kholinar fractions:
\begin{itemize}
	\item Helping \npcHighmarshal{} to defend the gates and secure the monastery
	\item Fighting alongside the \npcFractionCult{} against \npcHighmarshal{} and the parshmen
	\item Defending helpless parshmen and newly born Singers
\end{itemize}



\subsection{Event: The Thrill of Battle}

As the fights rise in the city, \npcThrill{}, the \npcThrillTitle{} gets attracted by the screams of battle and the blood of the wounded. \\

\eventComplication{The Thrill of Battle}{	
	\iconCompEvent
	\iconCompEvent
	\iconCompEvent
	\iconCompEvent\\ \vspace{0.15cm}
	\iconCompEvent
	\iconCompEvent
	\iconCompEvent
	\iconCompEvent\\ \vspace{0.15cm}
}

\subsubsection{Progressing the Event}

Each time a complication is rolled during a fight in Kholinar, you can spend it to progress the event by 1, trading an immediate Complication for the \npcThrill{}'s attention.

Each time a PC kills a foe in Kholinar while there are other enemies still alive, they must pass \skillCheck{15}{\skillDiscipline{}}. On a failure, progress the event by 1.


\subsubsection{Triggering the Event}

Once the final complication slot is filled, \npcThrill{} arrives to Kholinar. From now on, during all combat scenes, all characters (PCs and NPCs) are affected by the following rules:

\scriptedOption{Thrill of Murder}{
	Raise Stakes each time a melee attack is made. On a \iconComp{}, add +3 to that d20 roll and that attack grazes all characters adjacent to both the attacker and the defender (without spending focus).
}

\scriptedOption{Keep Fighting!}{
	If a character wishes to move away from the enemy or stop fighting (e.g. when there are no enemies left), they must make a \skillCheck{15}{\skillDiscipline}. On a failure they must keep fighting and attempt an attack against any character in the scene, including an ally. If the character passes the test, they no longer repeat it during next rounds of this combat scene.
}

\corruptedBox{
	Fighting in the battle influenced by \npcThrill{}, the \npcThrillTitle{}, can expedite the progression of an Enlightened Windrunner \seeCorruptedRules{}.
}



\subsection{Death Rattles}

When a nearby NPC dies, you can consider marking his death with a Death Rattle, especially if there is an Enlightened Truthwatcher in the party.

You can make up your own Death Rattles (either random nonsense or foreshadowing prophecy), or you can pick one from the examples presented in \nameref{chap:deathRattles}.

\corruptedBox{
	Hearing a Death Rattle sent by a \npcMoelach{} can expedite the progression of an Enlightened Truthwatcher \seeCorruptedRules{}.
}

\reasonBox{
	\npcMoelach{} will have its own arc dedicated for its influence in \nameref{chap:chapter7}. Do not expedite an Enlightened Truthwatcher more than once in this chapter.
}


\section{Proclamation}

Start the events of this chapter by reading the following:

\quoteBox{
	You see small crowd gathering around a woman wearing queen's colors who stands on some box in the middle of a street. She is accompanied by royal guards and is holding a piece of parchment. \\ \\
	The crowd around grows agitated with each passing moment.
}

Female herald is an \enemyExpert{} guarded by two \enemyQueensGuard{}s. The crowd throws rotten food at them, but they are ignoring it.

The characters can lead the riot against them or try to calm the people. If the fight starts, the herald flees, leaving her scroll behind for the PCs to find, but the guards will stand and fight against the people. If no fighting starts, after a few moments the herald reads the content of the scroll:

\quoteBox{
	We, \npcQueenName{}, by the grace of Allmighty and His Heralds, Queen of the Alethkar, do herby proclaim that all parsh people are the Voidbringers in hiding and therefore enemies to the realm and our faith. It is our order to exterminate all parsh people taking shelter behind the walls of Kholinar.
}

After reading the proclamation, the herald and the guards walk away, leaving angry and confused people behind.

If the riots start before the herald can read the order, she leaves the scroll behind in a hurry. Any character who can read in alethi will be able to learn its contents.

\spoilerBox{
	\npcQueen{} proclaiming parsmen to be Voidbringers and ordering their execution is referenced in \bookOB{}, Chapter 62: Research.
}

\section{Blood on the Streets}
\label{sec:bloodyStreets}

To reach the monastery, the characters will need to travel through streets of Kholinar. That will spawn multiple opportunities for smaller encounters on the streets. Use this section for some inspiration, especially if the characters take more time to reach the monastery.

However, before you allow the PCs to chart their course, read the following:

\quoteBox{
	When you step on the street, you see a smoke rising from the city walls. Some people run away from that direction. \\
	"The gate! They are attacking the gate!", one of them shouts.
}

From here, the characters can rush towards the city gates, sneak to the monastery or even attempt their luck to enter the palace. When they walk on the streets it is possible for the events below to happen:



\subsection{More Proclamations}
\label{ssec:bloodyStreetsProclamations}

If long enough time has passed since the previous queen’s proclamation read by her herald, a new one can be announced. Each proclamation starts with the same preamble:

\quoteBox{
	We, \npcQueenName{}, by the grace of Almighty and His Heralds, Queen of the Alethkar, do hereby proclaim that...
}

Each new proclamation errs on the edge of insanity, painting the picture of a queen descending into madness and religious fanaticism. None of the announcements seem to acknowledge ongoing riots. Examples may include:

\begin{itemize}
	\item All lights should be put out after dark and the spheres covered with cloth to make it harder for Voidbringers to find Kholinar in the dead of the night.
	
	\item People should pray to the Almighty out in the open each day, including highstorm days, to beg for absolution and mercy for Kholinar.
	
	\item Each household should sacrifice a chull to beg the Almighty for the King's victory over the Voidbringers on the Shattered Plains.
	
	\item All people, both men and women, should adjust their clothing to cover both freehand and safehand as a sign of modesty and humility.
	
	\item All people should face the Origin and sing ten songs to the ten heralds ten times a day.
	
	\item Tallew gathered on odd days of the week is proclaimed impure and should be burnt. Only tallew gathered on even days are admitted for consumption.
	
	\item People shall bring all of their foods and drinks to the ardents to be purified with prayers before they can be consumed.	
\end{itemize}



\subsection{Violent Clash}

When starting this encounter, read the following:

\quoteBox{
	You hear battle roars and cries of pain. The street is red with blood. You can't tell who started this fight, there are wounded on both sides.
}

Choose 2 groups of adversaries:
\begin{itemize}
	\item 1d4 \enemyWallGuard{}s
	\item 1d6 \enemyRioter{}s
	\item 1d6 \enemyCultist{}s (each \enemyCultistOfMoments{} of random kind)
\end{itemize}

Those groups fight with each other. They ignore the characters, but will attack them once they enter the conflict helping one of the sides.


\subsection{Helpless Parshmen}

When starting this encounter, read the following:

\quoteBox{
	You see a docile parshmen worker, hauling some boxes. Angry mob gathers around him. \\
	"Storming Voidbringer!", one of the rioters shouts as he raises his hammer. "Die, you filth!"
}

A \enemyPrshmen{} is surrounded by 1d6 \enemyRioter{}s. Unless the characters intervene, they will kill the worker. The party can also join them and execute queen's order.

You can repeat this encounter multiple times, changing some details of the scene:
\begin{itemize}
	\item You can change aggressors' fraction, cycling between \enemyRioter{}s, \enemyWallGuard{}s and \enemyCultist{}s.
	\item The parshmen can have different coloration of patterns and occupation.
	\item If the party is quick to follow queen's order, introduce moral dilemma by making subsequent parshmen even more helpless: make them old and barely moving, pregnant or just a kid.
\end{itemize}



\subsection{Hostile Rioters}

When you start this encounter, read the following:

\quoteBox{
	You did not see that coming. A heavy blow hits you in the back. "You are one of them, aren't you?", the rioter asks. "Where do you think you are going? This street is ours!"
}

The party is attacked by 1d6 + 2 \enemyRioter{}s. Each character who doesn't pass \skillCheck{15}{Perception} is Surprised in the first round of the fight. The enemies fight until there are only 2 of them left, at which moment they flee.



\subsection{Cult Procession}

1d4 + 4 \enemyCultist{}s (each \enemyCultistOfMoments{} of random kind) walk in procession through Kholinar streets. They seem oblivious to riots all around them and they will not attack if unprovoked. If the characters start the fight, they fight to the end. None of them carries anything valuable.

\spoilerBox{
	Similar procession is described in \bookOB{}, Chapter 61: Nightmare Made Manifest.
}




\subsection{Hungry Beggar}

When the party meets the beggar, read the following:

\quoteBox{
	You see an old and dirty man wearing rags, who begs on the side of the street. 
	“Please, I haven’t eaten for days”, he asks, reaching his only healthy hand towards you.
}

The characters can share the food with the beggar. It is a right thing to do, but they cannot expect anything in return, and food is scarce in the city...

You can repeat this encounter several times, changing the beggar NPC each time to someone new.



\subsection{Wounded Enemy}

When you introduce this encounter, read the following:

\quoteBox{
	You can see the rioters fought here not so long ago. The street is painted red with pools of blood, dead bodies lie scattered everywhere. You want to leave when you hear faint voice: "Please... help me..."
}

The party meets one badly wounded adversary (one from \enemyWallGuard{}, \enemyRioter{} and \enemyCultistOfMoments{}). They can help him if they so choose.

\reasonBox{
	To make this encounter more interesting, choose an NPC from the fraction the characters are the most hostile towards.
}



\subsection{Trouble at Home}

If the party decides to return to \npcGrandma{} during the riots, read the following:

\quoteBox{
	You hear grandma \npcGrandmaName{} even before you enter the building: \\ \\
	"Hey! Where are your manners! Just come here, you storming son of a chul! I will teach you! Hey, I am talking to you, you hear me? Come here! You will regret..."
}

Two \enemyRioter{}s and one \enemyWallGuard{} are looting the place. They ignore the old lady and the parshmen, but they will turn hostile when they notice the PCs entering. When only one of them is left standing, he will try to flee the building.



\subsection{It's a Hard Work}

The characters may decide that the riots are the best opportunity for some side job in the goods liberation and redistribution business. Depending on the value of the loot they are aiming for, consider no resistance, some competition in a form of a few other \enemyRioter{}s, hired guards protecting the place (use \enemyWallGuard{} stat block), or other adversaries from \bookWorldGuide{}.

\reasonBox{
	Try to ensure the party will not spend entire chapter just looting random places in the city. If you feel this activity is repeated too many times, let the spren urge them to focus on their goals, increase security of looted places, or, if that wasn't enough, trigger the events of \fullref{sec:enemyReborn}.
}
\section{The Palace}
\label{sec:closedPalace}

At this point in the story, there is little to be gained in the palace. If characters attempt to enter the palace, they will find closed and heavily guarded gates.

\quoteBox{
	Two knights guard the bared palace gates. Their armors shine in the sun, but there is a shadow falling over their eyes and the sense of wrongness to them. Their leader approaches you. "The palace is closed for the likes of you.", he says. "Go back where you belong!"
}

Two \enemyQueensGuard{}s keep watch on both sides of the gates. Their commander - \enemyQueensGuardLeader{} will not allow the characters to enter. If the characters stay here too long, they will be attacked by palace forces. When the fight starts, two more \enemyQueensGuard{} will reinforce palace forces after the first round of combat. All guards will fight to the end.

\corruptedBox{
	Approaching the gates of the Kholinar Palace influenced by \npcYeligNar{}, the \npcYeligNarTitle{}, can expedite an Enlightened Skybreaker \seeCorruptedRules{}.
}


\subsection{Talking Their Way In}

The characters may wish to persuade the guards to let them in - run it as a conversation.
\conversationDefence{\enemyQueensGuardLeader{}}{15}{18}{5}

His loyalty to the \npcQueen{} is incredible and he will positively respond to people with similar values.
All tests against him are made with disadvantage unless the character has very good reason to be allowed inside.

\conversationEnd{\enemyQueensGuardLeader{}}{1}



\subsection{Sneaking Inside}

The characters may decide to avoid the guards and sneak inside the palace behind their backs. You can run it as an endeavor.
\pursuitReq{6}{2}
All tests of the endeavor made against the palace guards are made with disadvantage as they are on high alert.

If the characters succeed, they can enter the palace without the guards knowing. Otherwise, they are attacked by them. They will fight to the death, but they will not pursue fleeing characters.



\subsection{Inside the Palace}

If the characters convinced the guards to let them in, they will be escorted by two \enemyQueensGuard{}s. They will be allowed to the throne room to swear devotion and loyalty to \npcQueen{} (or whatever was the reason the guards let them in). 

After that they will be escorted back to exit. Any sign of disobedience or lack of cooperation will turn the guards hostile and may result in lethal encounter.

\reasonBox{
	Do not allow \npcQueen{} to participate in a fight in this chapter - she will have her own boss fight in Act II.
}
\section{Battle at the Gates}
\label{sec:gatesBattle}

If the players managed to reach this location before completing \fullref{sec:monastery}, they will see \npcHighmarshal{} and her \enemyWallGuard{}s defending the gate from the \npcFractionCult{}. The PCs can join either side of the conflict. Start the encounter from  \fullref{ssec:gatesMain}.

If they decided to complete \fullref{sec:monastery} first, \npcHighmarshal{} successfully defended the gate by the time the characters reach it. In this case there is nothing more to do in here and you can skip the events of this section.

\roleplayBox{\npcHighmarshal{}}{
	\roleplayTips
		{Determined, inspiring, resourceful}
		{Protect the people of Kholinar}
		{Orange eyes and short, straight, black hair. Lean and grim, with couple scars on her face. Wears shining chainmail and an azure cloak.}
		{\npcHighmarshal{} didn't ask for that - she just was around when Kholinar needed a Wall Guard commander and she stepped up to take the responsibility. She is not even from Alethkar, but this does not matter - there are people who need her and she will defend them!}
}

\spoilerBox{
	\npcHighmarshal{} leading the defense of the city gates against the Cult of Moments is referenced in \bookOB{}, Chapter 70: Highmarshal Azure.
}



\subsection{Main Gate}
\label{ssec:gatesMain}

Start this encounter by reading the following:

\quoteBox{
	"Form the line! Shields high! Don't give them an inch!" - you hear orders shouted from a small group of defenders surrounded by an assault of cultists in white.
}

\enemyWallGuardHighmarshal{} leads a small squad of four \enemyWallGuard{}s against an assault of eight \enemyCultist{}s (each \enemyCultistOfMoments{} of a random kind). \npcHighmarshal{} will retreat once all her troops are dead, but all others will fight to the end.

From here, the characters can take the stair to join the fight in \fullref{ssec:gatesWalls}. They can also enter the gate tower to fight in \fullref{ssec:gatesBarracks}.

\npcHighmarshal{} and other combatants form this location will remain here, it is up to the players to drive off the enemies.




\subsection{Guard Barracks}
\label{ssec:gatesBarracks}

When the PCs enter this room, read the following:

\quoteBox{
	You enter the barrack room to see the last of the cultists being overwhelmed by the defenders.
}

Six \enemyWallGuard{}s are fighting here, and only two \enemyCultist{}s (random kinds) are remaining. All of them will fight to the end.

Both sides are vaguely familiar with the character's involvement during the fight in \fullref{ssec:gatesMain}, so the PCs cannot change sides now.

From here the characters can take an exit to reach \fullref{ssec:gatesWalls}.




\subsection{City Walls}
\label{ssec:gatesWalls}

When the PCs enter the walls, read the following:

\quoteBox{
	You enter just in time to see more cultists approaching on the walls. The lone guard will soon be overwhelmed...
}

Single \enemyWallGuard{} is about to be attacked by six \enemyCultist{}s (random kinds). All will fight to the end.

Both sides are vaguely familiar with the character's involvement during the fight in \fullref{ssec:gatesMain}, so the PCs cannot change sides now.

From here the characters can enter \fullref{ssec:gatesBarracks}.




\subsection{Aftermath}

The fighting stops once the PCs defeat all enemies in all three locations. The gates are now under control of the fraction the PCs were supporting during the fight.

\npcHighmarshal{} is gathering people to storm the Cultist Monastery (see \fullref{sec:monastery}). She invites the characters to join her.

If \npcHighmarshal{}'s forces were defeated by the PCs, she will retreat to recover and attack the Monastery alone.
\section{Ardent Monastery}
\label{sec:monastery}

Sooner or later, the characters are bound to attempt entering the Monastery. When they arrive to the location, they will start from \fullref{ssec:monasteryYard}.

\subsubsection{\npcHighmarshal{}'s Assault}

If the party was quick to get here, they will find mostly Ardents and \enemyCultist{}s here. If they stopped to help one of the sides during \fullref{sec:gatesBattle}, \npcHighmarshal{} will also storm the Monastery, and she will bring \enemyWallGuard{}s with her (if enough of them survived). 

The party will meet her in \fullref{ssec:monasterySoulcasters}, unless they are attacking the place together, allowing them to share most of the scenes.

\subsubsection{Past Actions}

Depending on players' actions so far, one of the sides of this conflict may be friendly towards the party. The PCs may storm the place alongside \npcHighmarshal{} or they may decide to help defend the place, hoping the newly-earned trust of the Ardents will help them achieve their goals.

If the characters were hostile to \enemyCultistOfMoments{}s earlier in the story, they might be recognized, turning peaceful encounter into hostile one.


\spoilerBox{
	\npcHighmarshal{} attacking the low monastery by the eastern gates is referenced in \bookOB{}, Chapter 73: Telling Which Stories.
}



\subsection{Monastery Yard}
\label{ssec:monasteryYard}

When the characters enter the Monastery, read the following:

\quoteBox{
	You are surrounded by serenity, as though this place was oblivious to the chaos raging behind its walls. You see two figures in white robes approaching you: \\ \\
	"Hello, child!", one of them says. "What are you looking for in this place?"
}

Two Ardents (use \enemyCultist{}'s stat blocks) guard the entrance. They are not hostile, unless they sense an incoming danger. If the party decides to attack them, two more \enemyCultist{}s will join the fight after the first round of combat.

The PCs may attempt to talk their way in. Run this as a conversation.
\conversationDefence{\enemyCultist}{11}{16}{3}
They live their lives devoted to the Almighty, and recently - the spren. They will respond positively to arguments suggesting the characters want to take part in the religious rituals.

\conversationEnd{\enemyCultist}{1}

If the players ask about the Unmade artifacts, the Ardents will deny existence of such items, calling it a heresy and a blasphemy. They will cease any further conversation and will turn hostile if the characters will press the topic.

From here, the characters can enter \fullref{ssec:monasteryHalls}.

\subsubsection{\npcHighmarshal{}'s Assault}

If \npcHighmarshal{} is storming the Monastery, there will be 4 \enemyCultist{}s and 3 \enemyWallGuard{}s fighting in here.




\subsection{Long Halls}
\label{ssec:monasteryHalls}

When the PCs enter the halls, read the following:

\quoteBox{
	Long, dark halls connect the whole complex. The walls of cold stone are decorated only with the sparse glow of torchlight. A few ardents walk in here.
}

Sometimes an Ardent or \enemyCultistOfMoments{} can be found here, but they will ignore the characters. The halls lead to \fullref{sec:monasteryChapel}. One of the closed doors hides \fullref{sec:monasteryFeast}. At the end of the corridors, there is \fullref{sec:monasteryStair}. The PCs can also reach \fullref{sec:monasteryCells} in the farthest parts of the Monastery.

\subsubsection{\npcHighmarshal{}'s Assault}

If \npcHighmarshal{} is storming the Monastery, there will be 2 \enemyCultist{}s and 2 \enemyWallGuard{}s fighting in here.



\subsection{Monastery Chapel}
\label{sec:monasteryChapel}

On the walls of the chapel there are paintings of spren, before them a group of kneeling people is praying - ardents in white robes lead the incantations, while the common people repeat after them:

\quoteBox{
	\centering
	"Oh, bringers of joy, remember us!" \\ "Remember us!" \\
	"Oh, dancers with flames, warm us!" \\ "Warm us!" \\
	"Oh, tenders of life, tend to us!" \\ "Tend to us!" \\
	"Oh, whispers of wind, whisper to us!" \\ "Whisper to us!" \\
}

Several \enemyCultist{}s (each \enemyCultistOfMoments{} of a different kind, no two repeating) pray in here with common people and brighteyes like. All of them ignore the PCs, but if they attack them, they will defend themselves. The \enemyCultist{}s will fight to the end, without any regard for the common people, who will attempt to run away.




\subsection{Lavish Feast}
\label{sec:monasteryFeast}

The door to this location is closed. If a character wishes to enter despite that, they will need to pass \skillCheck{15}{\skillStrength{} or \skillThievery{}}. In that case, read the following:

\quoteBox{
	You see a long table full of delights - meat, vegetables, cakes, exotic fruit - illuminated with spherelight from many chalices. More of the food lies on the floor and besides the walls. Some of it starts to rot. \\ \\
	A group of people in white robes feast on the delights. "You are not supposed to be here!", one of them shouts.
}

A \enemyRevelerDrinking{} and eight \enemyCultist{}s (of random kinds) feast in here. They will turn hostile unless a characters manages to quickly convince them to stand down - they would need to pass a \skillDeception{} or \skillPersuasion{} test against \enemyRevelerMasked{}'s Spiritual Defense with disadvantage.

If a combat starts, the enemies will fight to the end.

\subsubsection{\npcHighmarshal{}'s Assault}

If \npcHighmarshal{} is storming the Monastery, her forces will ignore this room, but some of \enemyWallGuard{}s can be called for help.

\subsubsection{Treasure}

In the chalices, there is 5d10 marks of spheres in total. \enemyRevelerDrinking{} carries a key to \fullref{sec:monasteryCells}.

The characters can also help themselves with the food or even bring something back home.

\corruptedBox{
	Participating in a lavish feast in the honor of \npcHeartRevel{}, the \npcHeartRevelTitle{}, may expedite an Enlightened Willshaper.
}



\newpage
\subsection{Locked Cells}
\label{sec:monasteryCells}

If the characters enter this area, read the following:

\quoteBox{
	This area of the Monastery seems abandoned. The halls are dark and nobody can be seen in here. \\ \\
	Suddenly, you hear a faint voice from one of the cells: "Please, have mercy... I am so hungry..."
}

Ardent \npcLockedArdent[s]{} is locked in one of the cells. She is imprisoned by \npcFractionCult{}. She hasn't seen anyone since she was put in here.

\subsubsection{Releasing \npcLockedArdent{}}

A character may open the jail using key found on one of the Revelers in \fullref{sec:monasteryFeast}. If they don't have it, they can attempt \skillCheck{20}{\skillThievery}.

\subsubsection{Talking with \npcLockedArdent{}}

\npcLockedArdent{} is too weak and hungry for a proper conversation. If fed, she will gladly answer any questions:

\scriptedOption{Who are you?}{
	I am \npcLockedArdent{}, an ardent from Devotary of Denial.
}

\scriptedOption{How long have you been locked in here?}{
	Hard to say, I lost track of time... What day do have...? Anyway, they came for me a day after \npcRebelArdent{} got assigned to Queen's Retinue... Did she told them something about me? I thought we were friends...
}

\scriptedOption{Can you help us navigate the Monastery?}{
	It isn't the one I've spent most of my life in, but it should be similar. I can be your guide, just please, don't leave me here!
}

\scriptedOption{Do you know about an Artifacts of the Unmade?}{
	The Voidbringers are destroyed, nothing left from them. There is no such item, and if there were, you wouldn't find them in here, this is a holy place. You are wasting your time... And please, don't pursue such unholy items, they will only curse your soul... Trust in the Almighty, He will guide your souls!
}




\subsection{Stairways}
\label{sec:monasteryStair}

When the PCs reach this place, read the following:

\quoteBox{
	You see a spiral stairs cut from the stone, descending into a darkness. Several people in white robes are standing nearby. "This area is for ardents only. Return to the chapel", one of them says.
}

Six \enemyCultist{}s (each \enemyCultistOfMoments{} of random kind) guard the stairway. They will turn hostile unless the PCs quickly go back. 

From here, the PCs can descend to \fullref{ssec:monasterySoulcasters}.

\subsubsection{\npcHighmarshal{}'s Assault}

If \npcHighmarshal{} is storming the Monastery, four \enemyWallGuard{} will fight the \enemyCultist{}s. The party can help in combat or sneak downstairs.




\subsection{\npcSoulcaster{} and Her Sister}
\label{ssec:monasterySoulcasters}

When the PCs reach this place, read the following:

\quoteBox{
	You see two female ardents sitting in the middle of the room. Green vines pierce through their pale skins, entangling their frail bodies. One of them wears a Soulcaster on her hand.
}

\npcSoulcaster{} and her sisters are \enemySoulcaster{}s, but only one of them has the Fabrial required for Soulcasting. They will not attack any aggressors - in a case of combat, they will attempt to flee or slow down their enemies. They will use Soulcaster as a weapon only when they will be left with no other choice.

No matter the cost, they will not part with the Fabrial. They are not part of the \npcFractionCult{} and they are not familiar with current events or political forces.

At the end of the room, the door leads to \fullref{ssec:monasteryReliquary}. Any character who bonded an Enlightened spren will feel a slight pull towards that area.

\spoilerBox{
	\npcSoulcaster{} is met in \bookOB{}, Chapter 81: Ithi and Her Sister.
}

\subsubsection{\npcHighmarshal{}'s Assault}

If \npcHighmarshal{} is storming the Monastery, she can be found in this room, attempting to persuade the Soulcaster sisters to join her and help the falling city.




\subsection{Forbidden Reliquary}
\label{ssec:monasteryReliquary}

When the PCs reach this area, read the following:

\quoteBox{
	The room is full of small chests, each standing on its own stone shelf and containing an ancient relic of ardents long gone. Two of the relics catch your attention - the artifacts that \npcSjaAnat{} described to you.
}

The party found \lootDustcaster{} and \lootGateOpener{}, securing their objective. Other items in this room have no practical use.

\reasonBox{
	As the characters get hold of the fabrials from the reliquary, don't reveal all their secrets just yet. Let there be some mystery around them. Don't worry, the book will provide explanations when they are needed.
}

\corruptedBox{
	Seeing \lootDustcaster{} for the first time, an artifact of \npcDustmother{}, the \npcDustmotherTitle{}, can expedite the progression of an Enlightened Dustbringer \seeCorruptedRules{}.
}

\corruptedBox{
	Seeing \lootGateOpener{} for the first time, an artifact of \npcBlackFisher{}, the \npcBlackFisherTitle{}, can expedite the progression of an Enlightened Elsecaller \seeCorruptedRules{}.
}




\subsection{Aftermath}

When the party secured the artifacts of the Unmade, trigger the events of \fullref{sec:enemyReborn} when they exit the monastery. The party will still be able to pursue other goals of this chapter, but the conditions in Kholinar will be much different.

\section{The Everstorm}
\label{sec:enemyReborn}

Introduce events of this section after the party completes \fullref{sec:monastery}.

\quoteBox{
	The sky darkens as the thick, black clouds gather from the west. The wind raises, blowing with unnatural speed. Soon after, a red lightning strikes a building. More of them follows soon after.\\ \\
	In the midst of this, the usually quiet parshmen fell to their knees, screaming a soul-piercing wails of unimaginable anguish. Their bones break, their skin tears to shreds as dozens of spikes puncture through their bodies. Soon after, they all lie in pools of their own blood, with enough life in them to suffer agony.
}

The party faces two\enemySingerWarform{}s and three \enemySingerNimbleform{}s. They are Prone, lying in blood, Stunned for two rounds and Surprised for one. Once they are back on their feet they will attack the characters or nearby NPCs.

If the party decides not to take shelter to hide from the Everstorm and hostile Singers you can run the following encounters:



\subsection{Ungrateful Servants}
\label{ssec:enemyRebornServants}

If the party returns to \npcGrandma{}'s home, read the following:

\quoteBox{
	The household you enter is unnaturally quiet. No shouting, no complaining, no insults or empty threats. The building sounds empty...
}

If the characters rushed here just after the Everstorm arrived, they will see \enemySingerWarform{} and \enemySingerNimbleform{} standing above terrified \npcGrandma{}, savoring their revenge and preparing for the kill. If someone enters the room, they will attack them instead.

If the characters are late, they will find \npcGrandma{} brutally murdered by her own servants. They joined other Singers spreading mayhem in the city - the house is empty now.



\subsection{Raging Everstorm}

You can run this encounter multiple times and/or in pair with other events.

\quoteBox{
	Gusts of wind grow heavier. Soon after, a series of lightning strikes from the sky, flashing red light on otherwise dark street.
}

Lightning bolts of Everstorm hit the area. For each character in the scene, roll as if \enemySingerStormform{} attacked them with Bolt of Lightning while having a disadvantage. If the attack misses the character, the bolt lands nearby, destroying the surrounding environment.



\subsection{Clash With the Singers}

Read the following:

\quoteBox{
	In the howling wind you could not hear them, before you walked just into them: a group of Singers, still covered in their own blood after the transformation, fighting a terrified group of people. Red flashes of lightning flare every few seconds, but in between everything is black.
}

A group of Singers (\enemySingerDireform{} or \enemySingerStormform{}, 1d4 \enemySingerWarform{}s and 1d4 \enemySingerNimbleform{}s) fight against some humans (2d6 \enemyCultist{}s or 2d6 \enemyRioter{}s or 2d4 \enemyWallGuard{}s). Both groups ignore the party, but the characters may join on any side if they wish. All involved will fight to the end.



\subsection{Massacre the Weak}

Read the following:

\quoteBox{
	You see a few Singers standing above a crying mother. One of them grabbed her infant and smashed it into the stone wall of a building. Another one lunged to the devastated woman soon after.
}

A \enemySingerDireform{} and two \enemySingerWarform{}s are committing the atrocities. They did not notice the party yet, allowing them to sneak away. If they fail to do so, the singers will turn hostile towards them soon after. They will fight to the end.


\subsection{Unfortunate Artifabrian}

You can start similar encounter when a PC uses a fabrial or a non-Enlightened Radiant surgebends.

\quoteBox{
	You see a group of guards fighting against two Singers with sparks of lightning dancing around their arms. One of the guards aims with her fabrial to one of them. Soon after a bright light flashes the street and the Singer falls to the ground, twitching. \\ \\
	
	Cheers don't last long - moments later a vicious, screaming spren darts from the above towards the fabrial. A majestic figure appeared shortly after. The singer rammed one of the guards with a lance and send another one flying in the air. With the fabrial secured, the figure flied away, its unnaturally long red robes flapping behind it on the wind.
}

\enemySingerHaevenlyOne{} ignores the characters as long as they are not using fabrials, but two wounded \enemySingerStormform{} (with roughly half of their Health) remain. Unless the characters flee right away, they will turn hostile towards them. 


\section{Dying Singer}

When characters complete their objective for this chapter, end it with following:

\quoteBox{
	You see a badly wounded Singer sitting on the ground. "I must reach the others, beyond the walls", she mumbles to herself.
}

The characters can help the wounded or finish her off. If they choose to be friendly, she will introduce herself as \npcSingerWouned{}. She will invite them to the camp next to the city, where all other Singers gather, and she will ask them to help her get there.

\subsection{Singer Camp}

If the characters travel to the Singer camp, read the following:

\quoteBox{
	You reach the camp that is setup in an alethi fashion. Two massive Singers stay on watch, accompanied with one sleek one standing behind. \\ \\
	"Stop! Who goes there?", one of them asks, raising his weapon.
}

Two \enemySingerWarform{}s and one \enemySingerEnvoyform{} named \npcSingerEnvoy{} stay at the gates. If the party wishes to enter the camp, they must convince them to let them through. Run this as a conversation encounter.

\conversationDefence{\npcSingerEnvoy{}}{15}{18}{6}
Her goal is to grow Singer's camp in strength and she will welcome every potential ally, but she is wary of traitorous humans. 

All tests done by the party have advantage if they are accompanied by a Singer. Also, if during \fullref{sec:bloodyStreets} the characters stood in the defense of parshmen, one of them will step from the camp and vouch for the party, counting as one successful test.

\conversationEnd{\npcSingerEnvoy{}}{2}
If the Singers fail to believe the characters, they will turn hostile at the end of the conversation.

If the party chooses violent solution, more Singers will soon join the fight. There is a full camp of them, the characters can't defeat them all and will soon need to flee back to the city.

\section{Journeying Onward}

Once you complete encounters with reborn Singers, advance characters to level 4 and proceed to \nameref{chap:chapter4}.